% SPDX-License-Identifier: GPL-3.0-or-later
% Copyright (C) 2026 Alessandro Greco
\chapter{Architettura della rete}\label{cap:06-architettura}
\aggiornatoal{23 settembre 2026}{tor 0.4.9.13; consenso pubblico del 23/09/2026 ore 15:00 UTC}

\begin{obiettivi}
Al termine di questo capitolo saprai:
\begin{itemize}[nosep]
  \item elencare i ruoli dei nodi della rete Tor e le responsabilità di ciascuno;
  \item descrivere i livelli del protocollo, dal canale TLS allo stream dell'applicazione;
  \item seguire, a grandi linee, tutto ciò che accade quando un utente apre una pagina web con Tor;
  \item capire come implementazioni e versioni diverse convivono nella stessa rete.
\end{itemize}
\end{obiettivi}

Prima di scendere nei dettagli dei capitoli successivi, guardiamo la rete dall'alto. Questo capitolo è una mappa: ogni elemento verrà ripreso e approfondito più avanti, e i rimandi indicano dove.

% ---------------------------------------------------------------------------
\section{I ruoli}\label{sec:arch-ruoli}

La rete Tor è formata da programmi che svolgono ruoli diversi, spesso più di uno contemporaneamente. La \cref{fig:arch-ruoli} li riassume.

\begin{figure}[htbp]
\centering
% SPDX-License-Identifier: GPL-3.0-or-later
% Copyright (C) 2026 Alessandro Greco
%% Testo alternativo per l'accessibilita' (usato solo nell'EPUB).
\altfig{Schema dei ruoli nella rete Tor. In alto, per raggiungere un sito web, il client attraversa guard, middle ed exit. Tre simboli rappresentano schematicamente le nove directory authority che votano e firmano il consenso, diffuso nella rete. Un bridge può fungere da primo nodo in caso di censura. In basso, per raggiungere un servizio onion senza relay di uscita, il client e il servizio stabiliscono ciascuno un circuito Tor verso il rendezvous point; i due collegamenti tratteggiati non mostrano tutti i relay intermedi.}
% Ruoli nella rete Tor
\begin{tikzpicture}[font=\sffamily\small,>=Stealth,
  aut/.style={relay,fill=torviola!15,draw=torviola,font=\sffamily\scriptsize},
  ex/.style={relay,fill=torverde!30}]
  \node[client] (c) at (0,0) {Client};
  \node[relay] (g) at (2.8,0) {Guard};
  \node[relay] (m) at (5.3,0) {Middle};
  \node[ex] (e) at (7.8,0) {Exit};
  \node[server] (s) at (10.6,0) {Sito web};
  \draw[cellaf] (c) -- (g); \draw[cellaf] (g) -- (m); \draw[cellaf] (m) -- (e); \draw[->,thick,torgrigio] (e) -- (s);
  \foreach \x/\n in {3.8/1,5.3/2,6.8/3} \node[aut] (a\n) at (\x,2.4) {aut.};
  \node[font=\sffamily\scriptsize,above=1pt of a2] {directory authority (9): consenso firmato};
  \draw[dotted,thick,torviola] (a2) -- node[font=\sffamily\scriptsize,right]{consenso} (m);
  \node[relay,fill=torgrigio!10,draw=torgrigio,font=\sffamily\scriptsize] (b) at (1.4,-2.2) {bridge};
  \node[font=\sffamily\scriptsize,align=center] at (1.4,-3.1) {primo salto al posto\\del guard, se censurati};
  \draw[dashed,torgrigio,->] (c) -- (b);
  \node[client,fill=torverde!15,draw=torverde!60!black] (o) at (8.8,-2.3) {Onion service};
  \node[relay,fill=torverde!25,font=\sffamily\scriptsize,align=center] (rp) at (5.3,-2.3) {rendez-\\vous};
  \draw[dashed,torverde!60!black,->] (o) -- node[font=\sffamily\scriptsize,below]{circuito} (rp);
  \draw[dashed,torverde!60!black,->] (c) to[bend right=18] node[font=\sffamily\scriptsize,above,pos=.68]{circuito client} (rp);
  \node[font=\sffamily\scriptsize,align=center] at (5.3,-3.35) {client e servizio si incontrano\\senza passare da un exit};
\end{tikzpicture}

\caption{I ruoli nella rete Tor. Il client costruisce un circuito attraverso guard, middle ed exit; le directory authority pubblicano il consenso che elenca i relay; i bridge sostituiscono il guard per chi è sotto censura; gli onion service si raggiungono incontrandosi in un punto di rendezvous, senza exit.}
\label{fig:arch-ruoli}
\end{figure}

\begin{description}
  \item[Client] Il programma dell'utente, per esempio quello incluso in Tor Browser. Sceglie i percorsi, costruisce i circuiti e offre alle applicazioni una porta SOCKS. Non accetta connessioni dalla rete e non si autentica (\cref{cap:07-canali}).
  \item[Relay] Un nodo che inoltra traffico per gli altri. A seconda dei flag assegnati dalle autorità e della sua configurazione può essere:
  \begin{itemize}
    \item \textbf{guard}: il primo salto dei circuiti di un client, scelto una volta e usato a lungo (\cref{sec:guard});
    \item \textbf{middle}: un salto intermedio, che non conosce né l'origine né la destinazione;
    \item \textbf{exit}: l'ultimo salto, che apre le connessioni verso Internet secondo la propria exit policy (\cref{cap:11-stream});
    \item \textbf{HSDir}: conserva i descrittori degli onion service (\cref{sec:os-hsdir});
    \item \textbf{directory cache}: ridistribuisce il consenso e i descrittori (\cref{cap:12-directory}).
  \end{itemize}
  \item[Directory authority] Nove server fidati che votano ogni ora sullo stato della rete e firmano il consenso; alcune misurano anche la banda dei relay.
  \item[Bridge] Un relay non pubblicato, usato come primo salto da chi non può raggiungere i relay pubblici, spesso dietro un pluggable transport (\cref{cap:16-censura}).
  \item[Onion service] Un servizio raggiungibile solo attraverso Tor, senza rivelare la propria posizione (\cref{cap:14-onion-services}).
\end{description}

Nel consenso pubblico del 23 settembre 2026 i relay sono 9358: 5306 con il flag \texttt{Guard}, 3177 con il flag \texttt{Exit} (1615 hanno entrambi). La banda dei guard vale il 70\,\% del totale, quella degli exit il 33\,\% (\cref{cap:13-percorso}).

% ---------------------------------------------------------------------------
\section{I livelli del protocollo}\label{sec:arch-livelli}

Il protocollo di Tor si può descrivere come una pila di livelli, ciascuno costruito sopra il precedente (\cref{fig:arch-livelli}).

\begin{figure}[htbp]
\centering
% SPDX-License-Identifier: GPL-3.0-or-later
% Copyright (C) 2026 Alessandro Greco
%% Testo alternativo per l'accessibilita' (usato solo nell'EPUB).
\altfig{Schema semplificato in sei livelli, dall'alto: applicazione HTTP su stream; comandi di stream BEGIN, DATA, END, SENDME (capitolo 11); messaggi RELAY con cifratura a strati (capitolo 10); gestione del circuito con CREATE2 e DESTROY come comandi di cella ed EXTEND2 come messaggio RELAY (capitolo 9); celle a lunghezza fissa di 514 byte nei link v4 o successivi, accanto a celle di lunghezza variabile (capitolo 8); canale TLS su TCP con CERTS e NETINFO (capitolo 7). Gli strati indicano funzioni, non una classificazione rigorosa di tipi di cella.}
% Pila dei livelli del protocollo
\begin{tikzpicture}[font=\sffamily\small,
  liv/.style={draw,rounded corners=2pt,minimum width=7cm,minimum height=0.8cm,align=center},
  rifcap/.style={font=\sffamily\scriptsize,text=torgrigio,anchor=west}]
  \node[liv,fill=torverde!20] (a) at (0,4) {Applicazione (HTTP, \ldots) su stream};
  \node[liv,fill=torverde!10] (s) at (0,3) {Stream: BEGIN, DATA, END, SENDME};
  \node[liv,fill=torviola!20] (r) at (0,2) {Messaggi relay cifrati a strati};
  \node[liv,fill=torviola!10] (c) at (0,1) {Circuiti: CREATE2, EXTEND2, DESTROY};
  \node[liv,fill=torgrigio!15] (l) at (0,0) {Celle fisse: 514 byte (v4+)};
  \node[liv,fill=torgrigio!5] (t) at (0,-1) {Canale: TLS su TCP, CERTS, NETINFO};
  \node[rifcap] at (3.7,3) {capitolo 11};
  \node[rifcap] at (3.7,2) {capitolo 10};
  \node[rifcap] at (3.7,1) {capitolo 9};
  \node[rifcap] at (3.7,0) {capitolo 8};
  \node[rifcap] at (3.7,-1) {capitolo 7};
\end{tikzpicture}

\caption{I livelli del protocollo di Tor e i capitoli che li descrivono. Ogni livello usa quello sottostante: gli stream viaggiano come messaggi relay, i messaggi relay dentro celle cifrate a strati, le celle dentro canali TLS.}
\label{fig:arch-livelli}
\end{figure}

\begin{description}
  \item[Canale] Una connessione TLS tra due nodi adiacenti, su cui si verifica l'identità del relay (\cref{cap:07-canali}).
  \item[Celle] Le unità di 514 byte che viaggiano sul canale (\cref{cap:08-celle}).
  \item[Circuito] Un percorso di più salti, costruito un salto alla volta con handshake che producono chiavi diverse per ogni salto (\cref{cap:09-circuiti}).
  \item[Cifratura a strati] I messaggi relay avvolti in uno strato per ogni salto (\cref{cap:10-cifratura-relay}).
  \item[Stream] Le connessioni TCP dell'applicazione, aperte dall'exit e controllate con finestre e controllo della congestione (\cref{cap:11-stream}).
\end{description}
A fianco di questa pila stanno il sistema delle directory, che fornisce ai client l'elenco dei relay, e il protocollo degli onion service, che usa i circuiti per incontrarsi in un punto di rendezvous.

% ---------------------------------------------------------------------------
\section{Una pagina web, dall'inizio alla fine}\label{sec:arch-percorso}

Seguiamo, a grandi linee, che cosa succede quando un utente avvia Tor Browser per la prima volta e apre una pagina.
\begin{enumerate}
  \item \textbf{Avvio.} tor contatta una directory authority o una delle circa duecento \emph{fallback directory} elencate nel codice e scarica il consenso, di cui verifica le firme; poi scarica i microdescrittori dei relay (\cref{cap:12-directory}).
  \item \textbf{Scelta dei guard.} Sceglie a caso, in proporzione alla banda, un piccolo campione di guard e ne usa due (\cref{sec:guard}).
  \item \textbf{Canale.} Apre una connessione TLS verso il guard e ne verifica l'identità con le celle \cella{CERTS} (\cref{cap:07-canali}).
  \item \textbf{Circuito.} Sceglie un exit adatto alla destinazione e un middle, e costruisce un circuito di tre salti con \cella{CREATE2} ed \texttt{EXTEND2} (\cref{cap:09-circuiti,cap:13-percorso}). In realtà tor costruisce alcuni circuiti in anticipo, per averli pronti quando servono.
  \item \textbf{Stream.} Il browser chiede, attraverso la porta SOCKS, una connessione verso \texttt{www.example.com:443}; tor invia un \texttt{BEGIN} all'exit, che risolve il nome e apre la connessione TCP (\cref{cap:11-stream}).
  \item \textbf{Dati.} Il browser parla HTTPS con il sito attraverso lo stream: l'exit vede il nome del sito e il traffico TLS, ma non chi lo richiede; il guard vede chi lo richiede, ma non che cosa; il middle non vede né l'uno né l'altro.
\end{enumerate}
L'intero processo di avvio richiede di solito pochi secondi. Molti dei meccanismi descritti nei capitoli successivi, come il padding, il controllo della congestione o i vanguards, lavorano in silenzio attorno a questa sequenza.

% ---------------------------------------------------------------------------
\section{Implementazioni e versioni}\label{sec:arch-versioni}

\subsection{Due implementazioni}

L'implementazione storica è \textbf{tor}, in C, che svolge tutti i ruoli: client, relay, autorità, onion service. Accanto a essa il Tor Project sviluppa \textbf{Arti}, una riscrittura in Rust (\cref{cap:30-arti}), oggi usata come client e per gli onion service. Nel consenso pubblico del 23 settembre 2026 nessun relay usa Arti: 6027 relay eseguono tor 0.4.9.12, 2520 la 0.4.9.11, e le versioni restanti sono per lo più precedenti della serie 0.4.9; 47 relay eseguono la versione di sviluppo 0.5.0.0-alpha-dev.

\subsection{Sottoprotocolli}

Come fa una rete di migliaia di nodi, aggiornati in tempi diversi, a introdurre nuove funzionalità? Tor non ragiona per numeri di versione del software, ma per \textbf{sottoprotocolli}: ogni funzionalità appartiene a un sottoprotocollo con un numero di versione, e ogni relay dichiara nel consenso quali versioni supporta, nella riga \texttt{pr}~\spec{tor-spec}{Subprotocol-based versioning}. Un relay della rete pubblica dichiara, per esempio:
\begin{lstlisting}[numbers=none]
pr Conflux=1 Cons=1-2 Desc=1-4 DirCache=2 FlowCtrl=1-2 HSDir=2 HSIntro=4-5
   HSRend=1-2 Link=3-5 LinkAuth=3 Microdesc=1-3 Padding=2 Relay=2-6
\end{lstlisting}
Un client che vuole usare ntor-v3 con un relay controlla che questo dichiari \texttt{Relay=4}; per CGO, \texttt{Relay=6}; per il controllo della congestione, \texttt{FlowCtrl=2}. Questi numeri sono \emph{capability} indipendenti, non versioni cumulative: \texttt{Relay=6} non implica formalmente \texttt{Relay=4} (\proposta{346}, che per questo preferisce chiamarle \enquote{capability di sottoprotocollo}), anche se in pratica tor le annuncia come un intervallo contiguo perché implementa le funzionalità in ordine. Il consenso elenca anche i sottoprotocolli \emph{richiesti}: una versione di tor che non li supporta sa di essere troppo vecchia e si rifiuta di funzionare, invece di mettere a rischio i propri utenti.

% ---------------------------------------------------------------------------
\section*{Riepilogo}
\begin{itemize}
  \item La rete è fatta di client, relay (guard, middle, exit, HSDir, directory cache), directory authority, bridge e onion service.
  \item Il protocollo è una pila: canale TLS, celle, circuiti, cifratura a strati, stream.
  \item Il guard conosce l'utente ma non la destinazione, l'exit la destinazione ma non l'utente.
  \item tor (C) e Arti (Rust) convivono; le funzionalità si negoziano per sottoprotocolli, dichiarati nella riga \texttt{pr} del consenso.
\end{itemize}

% ---------------------------------------------------------------------------
\section*{Esercizi}

\begin{esercizio}[1]\label{es:06-vista}
Per ciascun ruolo (guard, middle, exit, HSDir, directory authority) indica che cosa sa di un utente che visita un sito web attraverso Tor.
\end{esercizio}

\begin{esercizio}[1]\label{es:06-pr}
Nella riga \texttt{pr} riportata nel capitolo, quali funzionalità dei capitoli 9--11 supporta il relay? Potresti usarlo come exit con CGO?
\end{esercizio}

\begin{esercizio}[2]\label{es:06-livelli}
Una pagina web di 100\,KB viene scaricata attraverso Tor. Stima, livello per livello, quanti byte transitano sul collegamento tra client e guard, tenendo conto di celle, intestazioni dei messaggi relay e SENDME.
\end{esercizio}

% ---------------------------------------------------------------------------
\section*{Per approfondire}
\begin{itemize}
  \item R.~Dingledine, N.~Mathewson, P.~Syverson, \titolo{Tor: The Second-Generation Onion Router}~\cite{dingledine2004tor}: l'architettura originale, in gran parte ancora valida.
  \item Tor Metrics, \url{https://metrics.torproject.org/}: numeri aggiornati su relay, utenti e versioni.
\end{itemize}
